
\documentclass[fontset=none,11pt,a4paper]{ctexart}
\usepackage[margin=1.8cm]{geometry}
\usepackage{fontspec}
\usepackage{amssymb}
\usepackage{enumitem}
\usepackage{hyperref}
\usepackage{xcolor}
\usepackage{titlesec}
\usepackage{parskip}
\IfFontExistsTF{Songti SC}
  {\setCJKmainfont{Songti SC}\setCJKsansfont{Songti SC}\setCJKmonofont{Songti SC}}
  {\setCJKmainfont[BoldFont=FandolSong-Bold.otf]{FandolSong-Regular.otf}\setCJKsansfont[BoldFont=FandolSong-Bold.otf]{FandolSong-Regular.otf}\setCJKmonofont[BoldFont=FandolSong-Bold.otf]{FandolSong-Regular.otf}}
\IfFontExistsTF{TeX Gyre Termes}{\setmainfont{TeX Gyre Termes}}{\setmainfont{Times New Roman}}
\hypersetup{colorlinks=true, linkcolor=black, urlcolor=blue}
\setlength{\parindent}{0pt}
\setlist[itemize]{leftmargin=2em,itemsep=0.15em}
\setlist[enumerate]{leftmargin=2.6em,itemsep=0.15em}
\titleformat{\section}{\Large\bfseries}{\thesection.}{0.6em}{}
\titleformat{\subsection}{\large\bfseries}{\thesubsection}{0.6em}{}
\newcommand{\answerline}[1]{\underline{\hspace{#1}}}
\newcommand{\choicebox}[1]{$\square$~#1}
\newcommand{\blocktime}{\vspace{0.25em}\noindent 用时：\answerline{4cm}\par\vspace{0.6em}}
\newcommand{\tightpagebreak}{\vspace{0.4em}\hrule\vspace{0.6em}}

\title{\textbf{AIH 人机实验二：子 Agent 诊断性评测（T3）}}
\author{}
\date{}

\begin{document}
\maketitle
\vspace{-2em}

\textbf{作答说明：}本表单包含三部分：Block A 句子级诊断、Block B TimeBlock 时间推理、Block C 跨文本对齐。请只根据题面材料作答；不需要查阅整篇原文。Block A 的“时间信息原文”请填写原句中直接表达时间的短语，若没有则填“无”。

\section{Block A：句子级诊断（Agent 2--4）}
\blocktime
覆盖 Agent 2 时间信息原文抽取、Agent 3 下沉句判断、Agent 4 插叙/倒叙判断。

\paragraph{A-3-01 \quad 【53.9.4】}
\textbf{上下文（前五句 / 目标句 / 后五句）：}

\noindent 【53.8.4】高祖说：“好。”\par
\noindent 【53.8.5】于是便确定萧何为第一位，特恩许他带剑穿鞋上殿，上朝时可以不按礼仪小步快走。\par
\noindent 【53.9.1】高祖说：“我听说推荐贤才要受上等的奖赏。萧何的功劳虽然很高，经过鄂君的表彰就更加显赫了。”\par
\noindent 【53.9.2】于是根据鄂君原来受封的关内侯食邑，加封为安平侯。\par
\noindent 【53.9.3】当天，萧何父子兄弟十多人都封有食邑。\par
\noindent \textcolor{blue}{\textbf{【53.9.4】后又加封萧何两千户，这是因为高祖过去到咸阳服役时，萧何多送给自己二百钱的缘故。}}\par
\noindent 【53.10.1】汉十一年，陈豨反叛，高祖亲自率军到了邯郸。\par
\noindent 【53.10.2】平叛尚未结束，淮阴侯韩信又在关中谋反，吕后采用萧何的计策，杀了淮阴侯，此事记载在《淮阴侯列传》中。\par
\noindent 【53.10.3】高祖已经听说淮阴侯被杀，派遣使者拜丞相萧何为相国，加封五千户，并令五百名士卒、一名都尉做相国的卫队。\par
\noindent 【53.10.4】为此许多人都来祝贺，唯独召平表示哀悼。\par
\noindent 【53.10.5】召平原是秦朝的东陵侯。\par

\begin{itemize}
  \item 时间信息原文：\answerline{7.5cm}（无则填“无”）
  \item 是否为下沉句：\choicebox{是}\quad \choicebox{否}
  \item 是否为插叙/倒叙：\choicebox{是}\quad \choicebox{否}
\end{itemize}

\paragraph{A-3-02 \quad 【53.10.6】}
\textbf{上下文（前五句 / 目标句 / 后五句）：}

\noindent 【53.10.1】汉十一年，陈豨反叛，高祖亲自率军到了邯郸。\par
\noindent 【53.10.2】平叛尚未结束，淮阴侯韩信又在关中谋反，吕后采用萧何的计策，杀了淮阴侯，此事记载在《淮阴侯列传》中。\par
\noindent 【53.10.3】高祖已经听说淮阴侯被杀，派遣使者拜丞相萧何为相国，加封五千户，并令五百名士卒、一名都尉做相国的卫队。\par
\noindent 【53.10.4】为此许多人都来祝贺，唯独召平表示哀悼。\par
\noindent 【53.10.5】召平原是秦朝的东陵侯。\par
\noindent \textcolor{blue}{\textbf{【53.10.6】秦朝灭亡后，他沦为平民，家中贫穷，在长安城东种瓜。}}\par
\noindent 【53.10.7】他种的瓜味道甜美，所以社会上的人称它为“东陵瓜”，这是根据召平的封号来命名的。\par
\noindent 【53.10.8】召平对相国萧何说：“祸患从此开始了。皇上风吹日晒地统军在外，而您留守朝中，未遭战事之险，反而增加您的封邑并设置卫队，这是因为目前淮阴侯刚刚在京城谋反，对您的内心有所怀疑。设置卫队保护您，并非以此宠信您，希望您辞让封赏不受，把家产、资财全都捐助军队，那么皇上心里就会高兴。”\par
\noindent 【53.10.9】萧相国听从了他的计谋。\par
\noindent 【53.10.10】高帝果然非常欢喜。\par
\noindent 【53.11.1】汉十二年的秋天，黥布反叛，高祖亲自率军征讨他，多次派人来询问萧相国在做什么。\par

\begin{itemize}
  \item 时间信息原文：\answerline{7.5cm}（无则填“无”）
  \item 是否为下沉句：\choicebox{是}\quad \choicebox{否}
  \item 是否为插叙/倒叙：\choicebox{是}\quad \choicebox{否}
\end{itemize}

\paragraph{A-3-03 \quad 【8.26.1】}
\textbf{上下文（前五句 / 目标句 / 后五句）：}

\noindent 【8.24.2】秦地的百姓大失所望，然而心里恐惧，不敢不服从。\par
\noindent 【8.25.1】项羽派人回去报告楚怀王。\par
\noindent 【8.25.2】楚怀王说：“按照原来的约定办。”\par
\noindent 【8.25.3】项羽怨恨楚怀王不肯让他与沛公一起西进入关，而派他北上救赵，在天下诸侯争夺称王关中的约定中落在后面，他就说：“怀王这个人，我家项梁所立，没有什么功劳，凭什么主持约定。本来安定天下的，是诸位将领和我项籍。”\par
\noindent 【8.25.4】就假意推尊楚怀王为义帝，实际上不听从他的命令。\par
\noindent \textcolor{blue}{\textbf{【8.26.1】正月，项羽自立为西楚霸王，在梁、楚地区的九个郡称王，建都彭城。}}\par
\noindent 【8.26.2】背弃原来的约定，改立沛公为汉王，在巴、蜀、汉中称玉，建都南郑。\par
\noindent 【8.26.3】把关中瓜分为三，封立秦朝的三个将领：章邯为雍王，建都废丘：司马欣为塞王，建都栎阳；董翳为翟王，建都高奴。\par
\noindent 【8.26.4】封楚将瑕丘申阳为河南王，建都洛阳。\par
\noindent 【8.26.5】封赵将司马印为殷王，建都朝歌。\par
\noindent 【8.26.6】赵王歇迁徙代地称王。\par

\begin{itemize}
  \item 时间信息原文：\answerline{7.5cm}（无则填“无”）
  \item 是否为下沉句：\choicebox{是}\quad \choicebox{否}
  \item 是否为插叙/倒叙：\choicebox{是}\quad \choicebox{否}
\end{itemize}

\paragraph{A-3-04 \quad 【8.29.1】}
\textbf{上下文（前五句 / 目标句 / 后五句）：}

\noindent 【8.28.6】楚派萧公角攻打彭越，彭越大败萧公角。\par
\noindent 【8.28.7】陈余怨恨项羽不封自己为王，派夏说游说田荣，借兵攻打张耳。\par
\noindent 【8.28.8】齐借兵给陈余，击败了常山王张耳，张耳逃跑归附了汉王。\par
\noindent 【8.28.9】陈余从代接回赵王歇，又立为赵王，赵王就封陈余为代王。\par
\noindent 【8.28.10】项羽大怒，出兵北向击齐。\par
\noindent \textcolor{blue}{\textbf{【8.29.1】八月，汉王用韩信的计策，从故道回军，袭击雍王章邯。}}\par
\noindent 【8.29.2】章邯在陈仓迎击汉军，雍王兵败退走，在好畤停下来接战，又失败了，逃到废丘。\par
\noindent 【8.29.3】汉王随即平定了雍地。\par
\noindent 【8.29.4】向东到达咸阳，率军围困雍王于废丘，而派遣将领攻占了陇西、北地、上郡。\par
\noindent 【8.29.5】派将军薛欧、王吸出武关，借助王陵驻扎在南阳的兵力，迎接太公、吕后于沛县。\par
\noindent 【8.29.6】楚听到这一消息，出兵在阳夏阻挡，汉军不能前进。\par

\begin{itemize}
  \item 时间信息原文：\answerline{7.5cm}（无则填“无”）
  \item 是否为下沉句：\choicebox{是}\quad \choicebox{否}
  \item 是否为插叙/倒叙：\choicebox{是}\quad \choicebox{否}
\end{itemize}

\paragraph{A-3-05 \quad 【7.32.1】}
\textbf{上下文（前五句 / 目标句 / 后五句）：}

\noindent 【7.30.2】到了荥阳，各路败军都会合在一起，萧何也征发关中没有著籍的老弱全部来到荥阳，声势又振作起来。\par
\noindent 【7.30.3】楚军从彭城出发，常常乘胜追击败兵，与汉军在荥阳南面的京、索之间交战，汉军打败了楚军，楚军因此不能越过荥阳西进。\par
\noindent 【7.31.1】项王解救彭城，追赶汉王到达荥阳，田横乘机收复了齐国，立田荣子田广为齐王。\par
\noindent 【7.31.2】汉王在彭城战败，诸侯又都向楚背汉。\par
\noindent 【7.31.3】汉军驻扎在荥阳，修筑了一条甬道，与黄河相连，以便运取敖仓的粮食。\par
\noindent \textcolor{blue}{\textbf{【7.32.1】汉三年，项王屡次侵夺汉军的甬道，汉王粮食缺乏，恐慌起来，请求讲和，划分荥阳以西归汉。}}\par
\noindent 【7.32.2】项王想要答应他。\par
\noindent 【7.32.3】历阳侯范增说：“汉军容易对付，现在放掉他们，不予以消灭，以后一定要懊悔。”\par
\noindent 【7.32.4】项王就和范增加紧围攻荥阳。\par
\noindent 【7.32.5】汉王深为忧虑，就采用陈平的计策离间项王和范增。\par
\noindent 【7.32.6】项王的使者来了，给他准备了牛、羊、了齐全的丰盛筵席，打算端上去。\par

\begin{itemize}
  \item 时间信息原文：\answerline{7.5cm}（无则填“无”）
  \item 是否为下沉句：\choicebox{是}\quad \choicebox{否}
  \item 是否为插叙/倒叙：\choicebox{是}\quad \choicebox{否}
\end{itemize}

\paragraph{A-3-06 \quad 【53.11.1】}
\textbf{上下文（前五句 / 目标句 / 后五句）：}

\noindent 【53.10.6】秦朝灭亡后，他沦为平民，家中贫穷，在长安城东种瓜。\par
\noindent 【53.10.7】他种的瓜味道甜美，所以社会上的人称它为“东陵瓜”，这是根据召平的封号来命名的。\par
\noindent 【53.10.8】召平对相国萧何说：“祸患从此开始了。皇上风吹日晒地统军在外，而您留守朝中，未遭战事之险，反而增加您的封邑并设置卫队，这是因为目前淮阴侯刚刚在京城谋反，对您的内心有所怀疑。设置卫队保护您，并非以此宠信您，希望您辞让封赏不受，把家产、资财全都捐助军队，那么皇上心里就会高兴。”\par
\noindent 【53.10.9】萧相国听从了他的计谋。\par
\noindent 【53.10.10】高帝果然非常欢喜。\par
\noindent \textcolor{blue}{\textbf{【53.11.1】汉十二年的秋天，黥布反叛，高祖亲自率军征讨他，多次派人来询问萧相国在做什么。}}\par
\noindent 【53.11.2】萧相国因为皇上在军中，就在后方安抚勉励百姓，把自己的家财全都捐助军队，和讨伐陈豨时一样。\par
\noindent 【53.11.3】有一个门客劝告萧相国说：“您灭族的日子不远了。您位居相国，功劳数第一，还能够再加功吗？您当初进入关中就深得民心，至今十多年了，民众都亲附您，您还是那么勤勉地做事，与百姓关系和谐，受到爱戴。皇上之所以屡次询问您的情况，是害怕您震撼关中。如今您何不多买田地，采取低价、赊借等手段来败坏自己的声誉？这样，皇上的心才会安定。”\par
\noindent 【53.11.4】于是萧相国听从了他的计谋，高祖才非常高兴。\par
\noindent 【53.12.1】高祖征罢黥布军队回来，民众拦路上书，说相国低价强买百姓田地房屋数量极多。\par
\noindent 【53.12.2】高祖回到京城，相国进见。\par

\begin{itemize}
  \item 时间信息原文：\answerline{7.5cm}（无则填“无”）
  \item 是否为下沉句：\choicebox{是}\quad \choicebox{否}
  \item 是否为插叙/倒叙：\choicebox{是}\quad \choicebox{否}
\end{itemize}

\paragraph{A-3-07 \quad 【7.22.5】}
\textbf{上下文（前五句 / 目标句 / 后五句）：}

\noindent 【7.21.5】项王听到了这话，烹杀了劝说他的那个人。\par
\noindent 【7.22.1】项王派人向楚怀王报告请示。\par
\noindent 【7.22.2】楚怀王说：“按照约定办。”\par
\noindent 【7.22.3】项羽就尊楚怀王为义帝。\par
\noindent 【7.22.4】项王想自己称王，就先封诸侯将相为王。\par
\noindent \textcolor{blue}{\textbf{【7.22.5】对他们说：“天下最初发难的时候，暂时拥立诸侯后裔为王，以便讨伐秦朝。然而亲自身穿铠甲，手执兵器，率先起义，三年来风餐露宿，消灭秦朝，平定天下的，都是各位将相和我项籍的力量。只有义帝没有功劳，本来应该瓜分他的土地，封大家为王。”}}\par
\noindent 【7.22.6】将领们都说：“好。”\par
\noindent 【7.22.7】项王就分割天下，封将领们为侯王。\par
\noindent 【7.23.1】项王、范增疑心沛公将来占有天下。\par
\noindent 【7.23.2】但既已和解，又怕违背原约，诸侯反叛，他们就暗中商量说：“巴、蜀道路险恶，秦朝被迁徙的罪人都居住蜀地。”\par
\noindent 【7.23.3】于是就说：“巴、蜀也是关中地区。”\par

\begin{itemize}
  \item 时间信息原文：\answerline{7.5cm}（无则填“无”）
  \item 是否为下沉句：\choicebox{是}\quad \choicebox{否}
  \item 是否为插叙/倒叙：\choicebox{是}\quad \choicebox{否}
\end{itemize}

\paragraph{A-3-08 \quad 【8.23.11】}
\textbf{上下文（前五句 / 目标句 / 后五句）：}

\noindent 【8.23.6】沛公左司马曹无伤听说项王发怒，要攻打沛公，派人告诉项羽说：“沛公想要称王关中，令子婴为相，珍宝被他全部占有了。”\par
\noindent 【8.23.7】打算以此求得封赏。\par
\noindent 【8.23.8】亚父劝项羽进攻沛公。\par
\noindent 【8.23.9】当时项羽饱餐士卒，准备明日会战。\par
\noindent 【8.23.10】这时项羽兵四十万，号称百万。\par
\noindent \textcolor{blue}{\textbf{【8.23.11】沛公兵十万，号称二十万，兵力敌不过项羽。}}\par
\noindent 【8.23.12】恰巧项伯要救张良，夜间去见他。\par
\noindent 【8.23.13】用道理劝说项羽，项羽取消了进攻沛公的计划。\par
\noindent 【8.23.14】沛公带来了一百多骑兵，驰至鸿门，来见项羽，表示歉意。\par
\noindent 【8.23.15】项羽说：“这是你沛公左司马曹无伤向我说的。不然，我项羽何至于做这样的事。”\par
\noindent 【8.23.16】沛公因为樊哙、张良的缘故，得以脱身返回。\par

\begin{itemize}
  \item 时间信息原文：\answerline{7.5cm}（无则填“无”）
  \item 是否为下沉句：\choicebox{是}\quad \choicebox{否}
  \item 是否为插叙/倒叙：\choicebox{是}\quad \choicebox{否}
\end{itemize}


\newpage
\section{Block B：TimeBlock 推理（Agent 5--6）}
\blocktime
覆盖 Agent 5 时间标志物转换/补全，以及 Agent 6 TimeBlock 前后顺序判断。

\paragraph{B5-05 \quad Agent 5 / 时间标志物补全}
原文只写“八月”。结合前文“汉元年四月，在项羽麾下诸侯罢兵散归”的上下文，应补全为哪一项？


\textbf{材料：}【8.29.1】八月，汉王用韩信的计策，从故道回军，袭击雍王章邯。 【8.29.2】章邯在陈仓迎击汉军，雍王兵败退走，在好畤停下来接战，又失败了，逃到废丘。 【8.29.3】汉王随即平定了雍地。


\begin{enumerate}[label=\Alph*.]
  \item 汉元年四月
  \item 汉二年八月
  \item 秦二世三年八月
  \item 汉元年八月
\end{enumerate}
答案：\answerline{2.2cm}

\paragraph{B5-06 \quad Agent 5 / 时间标志物补全}
原文已经给出“汉十二年的秋天”。以下哪一项最适合作为该 TimeBlock 的时间标志物？


\textbf{材料：}【53.11.1】汉十二年的秋天，黥布反叛，高祖亲自率军征讨他，多次派人来询问萧相国在做什么。 【53.11.2】萧相国因为皇上在军中，就在后方安抚勉励百姓，把自己的家财全都捐助军队，和讨伐陈豨时一样。


\begin{enumerate}[label=\Alph*.]
  \item 汉十二年秋
  \item 汉十一年秋
  \item 汉三年秋
  \item 秦二世三年秋
\end{enumerate}
答案：\answerline{2.2cm}

\paragraph{B6-05 \quad Agent 6 / TimeBlock 顺序判断}
判断 A、B 两个 TimeBlock 的时间先后。


\textbf{材料 A：}【53.7.1】汉五年，已经消灭了项羽，平定了天下，于是论功行赏。

\textbf{材料 B：}【53.5.2】汉二年，汉王与各路诸侯攻打楚军，萧何守卫关中，侍奉太子，治理栎阳。 【53.5.3】制定法令、规章，建立宗庙、社稷、宫室、县邑，萧何总是禀报汉王，得到汉王同意，准许施行这些政事；如果来不及禀报汉王，有些事就酌情处理，等汉王回来再向他汇报。


\begin{enumerate}[label=\Alph*.]
  \item A 早于 B
  \item B 早于 A
  \item A 与 B 为同一事件或大致同时
  \item 仅凭材料无法判断
\end{enumerate}
答案：\answerline{2.2cm}

\paragraph{B6-06 \quad Agent 6 / TimeBlock 顺序判断}
判断 A、B 两段是否为同一历史事件或大致同时发生。


\textbf{材料 A：}【8.23.6】沛公左司马曹无伤听说项王发怒，要攻打沛公，派人告诉项羽说：“沛公想要称王关中，令子婴为相，珍宝被他全部占有了。”

\textbf{材料 B：}【7.17.6】沛公左司马曹无伤派人对项羽说：“沛公想称王关中，使子婴为相，占有了全部珍宝。”


\begin{enumerate}[label=\Alph*.]
  \item A 早于 B
  \item B 早于 A
  \item A 与 B 为同一事件或大致同时
  \item 仅凭材料无法判断
\end{enumerate}
答案：\answerline{2.2cm}


\newpage
\section{Block C：跨文本对齐（Agent 9）}
\blocktime
请选择最可能描述同一历史过程的 source。每题只有一个最佳答案。

\paragraph{C9-09 \quad Agent 9 / 跨文本对齐}
哪一个候选 source 最可能与 target 描述同一历史过程，从而可以向 target 传播时间信息？

\textbf{Target：}\textcolor{blue}{\textbf{【53.4.1】等到刘邦起事做了沛公，萧何常常做为他的助手督办公务。}}

\begin{enumerate}[label=\Alph*.]
  \item
\noindent 【8.21.7】乘胜追击，彻底打垮了秦军。\par
\noindent 【8.22.1】汉元年十月，沛公的军队先于各路诸侯到达霸上。\par
\noindent 【8.22.2】秦王于婴素车白马，用丝带系着脖子，封了皇帝的印玺和符节，在轵道旁投降。\par
\noindent 【8.22.3】将领们有的主张杀死秦王。\par
\noindent 【8.22.4】沛公说：“当初楚怀王派遣我，本来是因为我能宽大容人。况且人家已经降服，又杀死人家，不吉利。”\par
\noindent 【8.22.5】于是就把秦王交给了官吏，向西进入咸阳。\par
\noindent 【8.22.6】沛公想要留在宫殿中休息，樊哙、张良劝说后，才封闭了秦宫的贵重珍宝、财物和库房，回军霸上。\par
\noindent 【8.22.7】召集各县的父老、豪杰说：“父老们苦于秦朝的严刑峻法已经很久了，诽谤朝政的要灭族，相聚议论的要在街市上处斩。我和诸侯们约定，先入关的在关中称王，我应当称王关中。同父老们约定，法律只有三章：杀人的处死，伤人和抢劫的处以与所犯罪相当的刑罚。其余的秦朝法律全部废除。官吏和百姓都要安居如故。我所以到这里来，是为父老们除害，不会有欺凌暴虐的行为，不要害怕。我所以回军霸上，是等待诸侯们到来制定共同遵守的纪律。”\par
\noindent 【8.22.8】沛公派人与秦朝官吏巡行县城乡间，告谕百姓。\par
\noindent 【8.22.9】秦地的百姓大为高兴，争先恐后地拿出牛羊酒食款待士兵。\par
\noindent 【8.22.10】沛公又谦让不肯接受，说：“仓库的谷子很多，不缺乏，不愿破费百姓。”百姓更加高兴，唯恐沛公不做秦王。\par
  \item
\noindent 【8.25.2】楚怀王说：“按照原来的约定办。”\par
\noindent 【8.25.3】项羽怨恨楚怀王不肯让他与沛公一起西进入关，而派他北上救赵，\par
\noindent 【8.25.4】在天下诸侯争夺称王关中的约定中落在后面，他就说：“怀王这个人，我家项梁所立，没有什么功劳，凭什么主持约定。本来安定天下的，是诸位将领和我项籍。”\par
\noindent 【8.25.5】就假意推尊楚怀王为义帝，实际上不听从他的命令。\par
\noindent 【8.26.1】正月，项羽自立为西楚霸王，在梁、楚地区的九个郡称王，建都彭城。\par
\noindent 【8.26.2】背弃原来的约定，改立沛公为汉王，在巴、蜀、汉中称玉，建都南郑。\par
\noindent 【8.26.3】把关中瓜分为三，封立秦朝的三个将领：章邯为雍王，建都废丘：司马欣为塞王，建都栎阳；董翳为翟王，建都高奴。\par
\noindent 【8.26.4】封楚将瑕丘申阳为河南王，建都洛阳。\par
\noindent 【8.26.5】封赵将司马印为殷王，建都朝歌。\par
\noindent 【8.26.6】赵王歇迁徙代地称王。\par
\noindent 【8.26.7】封赵将张耳为常山王，建都襄国。\par
  \item
\noindent 【8.10.11】在沛县衙门的庭院里祭祀黄帝和蚩尤，又用牲血衅鼓旗。\par
\noindent 【8.10.12】旗子一律红色，因为刘季所杀蛇是白帝的儿子，杀蛇的是赤帝的儿子，所以崇尚赤色。\par
\noindent 【8.10.13】于是少年子弟和有势的官吏，如萧何、曹参、樊哙等人，都为沛公征集兵员，集合了两三千人，攻打胡陵、方与，回军固守丰邑。\par
\noindent 【8.11.1】秦二世二年，陈胜将领周章的军队西至戏水而还。\par
\noindent 【8.11.2】燕、赵、齐、魏都自立为王。\par
\noindent 【8.11.3】项梁、项羽起兵于吴。\par
\noindent 【8.11.4】秦泗水郡郡监平率兵围丰，两天后，沛公出兵应战，打败了秦军。\par
\noindent 【8.11.5】沛公命令雍齿守卫丰邑，自己引兵赴薛。\par
\noindent 【8.11.6】泗水郡郡守壮在薛战败，逃到戚。\par
\noindent 【8.11.7】沛公左司马擒获泗水郡郡守壮，杀死了他。\par
\noindent 【8.11.8】沛公回军亢父，到了方与，没有交战。\par
  \item
\noindent 【8.10.4】刘季用帛写了一封信，\par
\noindent 【8.10.5】射到城上，告诉沛县父老说：“天下苦于秦朝的暴政已经很久了。现在父老为沛令守城，但各国诸侯都已起事，（一旦城破，）就要屠戮沛县。如果沛县父老共同起来杀死沛令，选择子弟中可以立为首领的做领导，以响应诸侯军，那就能保全身家性命。不然的话，父子全遭杀害，死得毫无意义。”\par
\noindent 【8.10.6】父老们就率领子弟共同杀了沛令，打开城门，迎接刘季，想让他做柿县县令。\par
\noindent 【8.10.7】刘季说：“天下正在混乱当中，诸侯都已起事，如果推选的将领不胜任，就会一败涂地。我不是吝惜自己的生命，只怕才劣力薄，不能保全父兄子弟。这是件大事，希望另外共同推选一位能够胜任的人。”\par
\noindent 【8.10.8】萧何、曹参等都是文官，看重身家性命，怕事情不成，秦朝会诛灭他们的全族，所以都推让刘季。\par
\noindent 【8.10.9】父老们都说：“我们平时听到刘季许多奇异的事情，看来刘季是该显贵的。而且又经过占卜，没有比刘季更吉利的。”\par
\noindent 【8.10.10】这时刘季再三谦让，大家都不敢担任，最后还是立刘季为沛公。\par
\noindent 【8.10.11】在沛县衙门的庭院里祭祀黄帝和蚩尤，又用牲血衅鼓旗。\par
\noindent 【8.10.12】旗子一律红色，因为刘季所杀蛇是白帝的儿子，杀蛇的是赤帝的儿子，所以崇尚赤色。\par
\noindent 【8.10.13】于是少年子弟和有势的官吏，如萧何、曹参、樊哙等人，都为沛公征集兵员，集合了两三千人，攻打胡陵、方与，回军固守丰邑。\par
\noindent 【8.11.1】秦二世二年，陈胜将领周章的军队西至戏水而还。\par
\end{enumerate}
答案：\answerline{2.2cm}

\paragraph{C9-10 \quad Agent 9 / 跨文本对齐}
哪一个候选 source 最可能与 target 描述同一历史过程，从而可以向 target 传播时间信息？

\textbf{Target：}\textcolor{blue}{\textbf{【53.5.1】汉王领兵东进，平定三秦，萧何以丞相的身分留守治理巴蜀，安抚民众，发布政令，供给军队粮草。}}

\begin{enumerate}[label=\Alph*.]
  \item
\noindent 【8.25.2】楚怀王说：“按照原来的约定办。”\par
\noindent 【8.25.3】项羽怨恨楚怀王不肯让他与沛公一起西进入关，而派他北上救赵，\par
\noindent 【8.25.4】在天下诸侯争夺称王关中的约定中落在后面，他就说：“怀王这个人，我家项梁所立，没有什么功劳，凭什么主持约定。本来安定天下的，是诸位将领和我项籍。”\par
\noindent 【8.25.5】就假意推尊楚怀王为义帝，实际上不听从他的命令。\par
\noindent 【8.26.1】正月，项羽自立为西楚霸王，在梁、楚地区的九个郡称王，建都彭城。\par
\noindent 【8.26.2】背弃原来的约定，改立沛公为汉王，在巴、蜀、汉中称玉，建都南郑。\par
\noindent 【8.26.3】把关中瓜分为三，封立秦朝的三个将领：章邯为雍王，建都废丘：司马欣为塞王，建都栎阳；董翳为翟王，建都高奴。\par
\noindent 【8.26.4】封楚将瑕丘申阳为河南王，建都洛阳。\par
\noindent 【8.26.5】封赵将司马印为殷王，建都朝歌。\par
\noindent 【8.26.6】赵王歇迁徙代地称王。\par
\noindent 【8.26.7】封赵将张耳为常山王，建都襄国。\par
  \item
\noindent 【8.21.7】乘胜追击，彻底打垮了秦军。\par
\noindent 【8.22.1】汉元年十月，沛公的军队先于各路诸侯到达霸上。\par
\noindent 【8.22.2】秦王于婴素车白马，用丝带系着脖子，封了皇帝的印玺和符节，在轵道旁投降。\par
\noindent 【8.22.3】将领们有的主张杀死秦王。\par
\noindent 【8.22.4】沛公说：“当初楚怀王派遣我，本来是因为我能宽大容人。况且人家已经降服，又杀死人家，不吉利。”\par
\noindent 【8.22.5】于是就把秦王交给了官吏，向西进入咸阳。\par
\noindent 【8.22.6】沛公想要留在宫殿中休息，樊哙、张良劝说后，才封闭了秦宫的贵重珍宝、财物和库房，回军霸上。\par
\noindent 【8.22.7】召集各县的父老、豪杰说：“父老们苦于秦朝的严刑峻法已经很久了，诽谤朝政的要灭族，相聚议论的要在街市上处斩。我和诸侯们约定，先入关的在关中称王，我应当称王关中。同父老们约定，法律只有三章：杀人的处死，伤人和抢劫的处以与所犯罪相当的刑罚。其余的秦朝法律全部废除。官吏和百姓都要安居如故。我所以到这里来，是为父老们除害，不会有欺凌暴虐的行为，不要害怕。我所以回军霸上，是等待诸侯们到来制定共同遵守的纪律。”\par
\noindent 【8.22.8】沛公派人与秦朝官吏巡行县城乡间，告谕百姓。\par
\noindent 【8.22.9】秦地的百姓大为高兴，争先恐后地拿出牛羊酒食款待士兵。\par
\noindent 【8.22.10】沛公又谦让不肯接受，说：“仓库的谷子很多，不缺乏，不愿破费百姓。”百姓更加高兴，唯恐沛公不做秦王。\par
  \item
\noindent 【8.28.7】齐借兵给陈余，击败了常山王张耳，张耳逃跑归附了汉王。\par
\noindent 【8.28.8】陈余从代接回赵王歇，又立为赵王，赵王就封陈余为代王。\par
\noindent 【8.28.9】项羽大怒，出兵北向击齐。\par
\noindent 【8.29.1】八月，汉王用韩信的计策，从故道回军，袭击雍王章邯。\par
\noindent 【8.29.2】章邯在陈仓迎击汉军，雍王兵败退走，在好畤停下来接战，又失败了，逃到废丘。\par
\noindent 【8.29.3】汉王随即平定了雍地。\par
\noindent 【8.29.4】向东到达咸阳，率军围困雍王于废丘，而派遣将领攻占了陇西、北地、上郡。\par
\noindent 【8.29.5】派将军薛欧、王吸出武关，借助王陵驻扎在南阳的兵力，迎接太公、吕后于沛县。\par
\noindent 【8.29.6】楚听到这一消息，出兵在阳夏阻挡，汉军不能前进。\par
\noindent 【8.29.7】楚让原吴县县令郑昌为韩王，抵抗汉军。\par
\noindent 【8.30.1】二年，汉王东出略取城邑，塞王司马欣、翟王董翳、河南王申阳都投降了。\par
  \item
\noindent 【8.14.4】沛公和项羽正在攻打陈留，听说项梁死了，带兵和吕将军一起向东进发。\par
\noindent 【8.14.5】吕臣驻扎在彭城东面，项羽驻扎在彭城西面，沛公驻扎在肠。\par
\noindent 【8.15.1】章邯已经打垮了项梁的军队，以为楚地的敌人不用担心了，就渡过黄河，北进攻打赵地，大破赵军。\par
\noindent 【8.15.2】这个时候，赵歇为赵王，秦将王离围困赵歇于巨鹿城。\par
\noindent 【8.15.3】（被围在巨鹿的军队，）这就是所谓的“河北之军”。\par
\noindent 【8.16.1】秦二世三年，楚怀王看到项梁的军队被打垮了，心里恐惧，迁离盱台，建都彭城，合并吕臣、项羽的军队，亲自统率。\par
\noindent 【8.16.2】以沛公任砀郡长，封为武安侯，统领砀郡的军队。\par
\noindent 【8.16.3】封项羽为长安侯，号为鲁公。\par
\noindent 【8.16.4】吕臣任司徒，他的父亲吕青作令尹。\par
\noindent 【8.17.1】赵多次请求救援，楚怀王就以宋义为上将军，项羽为次将，范增为未将，北上救赵。\par
\noindent 【8.17.2】命令沛公西出略地，打入关中。\par
\end{enumerate}
答案：\answerline{2.2cm}

\paragraph{C9-11 \quad Agent 9 / 跨文本对齐}
哪一个候选 source 最可能与 target 描述同一历史过程，从而可以向 target 传播时间信息？

\textbf{Target：}\textcolor{blue}{\textbf{【53.10.2】平叛尚未结束，淮阴侯韩信又在关中谋反，吕后采用萧何的计策，杀了淮阴侯，此事记载在《淮阴侯列传》中。}}

\begin{enumerate}[label=\Alph*.]
  \item
\noindent 【8.25.1】项羽派人回去报告楚怀王。\par
\noindent 【8.25.2】楚怀王说：“按照原来的约定办。”\par
\noindent 【8.25.3】项羽怨恨楚怀王不肯让他与沛公一起西进入关，而派他北上救赵，\par
\noindent 【8.25.4】在天下诸侯争夺称王关中的约定中落在后面，他就说：“怀王这个人，我家项梁所立，没有什么功劳，凭什么主持约定。本来安定天下的，是诸位将领和我项籍。”\par
\noindent 【8.25.5】就假意推尊楚怀王为义帝，实际上不听从他的命令。\par
\noindent 【8.26.1】正月，项羽自立为西楚霸王，在梁、楚地区的九个郡称王，建都彭城。\par
\noindent 【8.26.2】背弃原来的约定，改立沛公为汉王，在巴、蜀、汉中称玉，建都南郑。\par
\noindent 【8.26.3】把关中瓜分为三，封立秦朝的三个将领：章邯为雍王，建都废丘：司马欣为塞王，建都栎阳；董翳为翟王，建都高奴。\par
\noindent 【8.26.4】封楚将瑕丘申阳为河南王，建都洛阳。\par
\noindent 【8.26.5】封赵将司马印为殷王，建都朝歌。\par
\noindent 【8.26.6】赵王歇迁徙代地称王。\par
  \item
\noindent 【8.20.13】回军攻打胡阳，遇到番君的别将梅鋗，与他一起，迫使析县、郦县投降。\par
\noindent 【8.20.14】派遣魏人宁昌出使秦关中，使者没有回来。\par
\noindent 【8.20.15】这时章邯已经带领全军在赵地投降项羽了。\par
\noindent 【8.21.1】起初，项羽和宋义北进援救赵，等到项羽杀死宋义，代替他为上将军，许多将领和黥布都从属项羽。\par
\noindent 【8.21.2】打垮了秦将王离的军队，使章邯投降，诸侯都归附了他。\par
\noindent 【8.21.3】等到赵高已经杀了秦二世，派人来见沛公，想要定约瓜分关中称王，沛公以为是诈骗，就采用张良的计策，派郦生、陆贾去游说秦军将领，用私利相诱，趁机袭击武关，攻破了关口。\par
\noindent 【8.21.4】又和秦军在蓝田南面交战，增设疑兵，多树旗帜，所经过的地方不许掳掠。\par
\noindent 【8.21.5】秦地的群众很高兴，秦军懈怠了，因此大破秦军。\par
\noindent 【8.21.6】又在蓝田北面接战，再次打败秦军。\par
\noindent 【8.21.7】乘胜追击，彻底打垮了秦军。\par
\noindent 【8.22.1】汉元年十月，沛公的军队先于各路诸侯到达霸上。\par
  \item
\noindent 【8.28.7】齐借兵给陈余，击败了常山王张耳，张耳逃跑归附了汉王。\par
\noindent 【8.28.8】陈余从代接回赵王歇，又立为赵王，赵王就封陈余为代王。\par
\noindent 【8.28.9】项羽大怒，出兵北向击齐。\par
\noindent 【8.29.1】八月，汉王用韩信的计策，从故道回军，袭击雍王章邯。\par
\noindent 【8.29.2】章邯在陈仓迎击汉军，雍王兵败退走，在好畤停下来接战，又失败了，逃到废丘。\par
\noindent 【8.29.3】汉王随即平定了雍地。\par
\noindent 【8.29.4】向东到达咸阳，率军围困雍王于废丘，而派遣将领攻占了陇西、北地、上郡。\par
\noindent 【8.29.5】派将军薛欧、王吸出武关，借助王陵驻扎在南阳的兵力，迎接太公、吕后于沛县。\par
\noindent 【8.29.6】楚听到这一消息，出兵在阳夏阻挡，汉军不能前进。\par
\noindent 【8.29.7】楚让原吴县县令郑昌为韩王，抵抗汉军。\par
\noindent 【8.30.1】二年，汉王东出略取城邑，塞王司马欣、翟王董翳、河南王申阳都投降了。\par
  \item
\noindent 【8.10.4】刘季用帛写了一封信，\par
\noindent 【8.10.5】射到城上，告诉沛县父老说：“天下苦于秦朝的暴政已经很久了。现在父老为沛令守城，但各国诸侯都已起事，（一旦城破，）就要屠戮沛县。如果沛县父老共同起来杀死沛令，选择子弟中可以立为首领的做领导，以响应诸侯军，那就能保全身家性命。不然的话，父子全遭杀害，死得毫无意义。”\par
\noindent 【8.10.6】父老们就率领子弟共同杀了沛令，打开城门，迎接刘季，想让他做柿县县令。\par
\noindent 【8.10.7】刘季说：“天下正在混乱当中，诸侯都已起事，如果推选的将领不胜任，就会一败涂地。我不是吝惜自己的生命，只怕才劣力薄，不能保全父兄子弟。这是件大事，希望另外共同推选一位能够胜任的人。”\par
\noindent 【8.10.8】萧何、曹参等都是文官，看重身家性命，怕事情不成，秦朝会诛灭他们的全族，所以都推让刘季。\par
\noindent 【8.10.9】父老们都说：“我们平时听到刘季许多奇异的事情，看来刘季是该显贵的。而且又经过占卜，没有比刘季更吉利的。”\par
\noindent 【8.10.10】这时刘季再三谦让，大家都不敢担任，最后还是立刘季为沛公。\par
\noindent 【8.10.11】在沛县衙门的庭院里祭祀黄帝和蚩尤，又用牲血衅鼓旗。\par
\noindent 【8.10.12】旗子一律红色，因为刘季所杀蛇是白帝的儿子，杀蛇的是赤帝的儿子，所以崇尚赤色。\par
\noindent 【8.10.13】于是少年子弟和有势的官吏，如萧何、曹参、樊哙等人，都为沛公征集兵员，集合了两三千人，攻打胡陵、方与，回军固守丰邑。\par
\noindent 【8.11.1】秦二世二年，陈胜将领周章的军队西至戏水而还。\par
\end{enumerate}
答案：\answerline{2.2cm}

\paragraph{C9-12 \quad Agent 9 / 跨文本对齐}
哪一个候选 source 最可能与 target 描述同一历史过程，从而可以向 target 传播时间信息？

\textbf{Target：}\textcolor{blue}{\textbf{【53.10.6】秦朝灭亡后，他沦为平民，家中贫穷，在长安城东种瓜。}}

\begin{enumerate}[label=\Alph*.]
  \item
\noindent 【8.25.1】项羽派人回去报告楚怀王。\par
\noindent 【8.25.2】楚怀王说：“按照原来的约定办。”\par
\noindent 【8.25.3】项羽怨恨楚怀王不肯让他与沛公一起西进入关，而派他北上救赵，\par
\noindent 【8.25.4】在天下诸侯争夺称王关中的约定中落在后面，他就说：“怀王这个人，我家项梁所立，没有什么功劳，凭什么主持约定。本来安定天下的，是诸位将领和我项籍。”\par
\noindent 【8.25.5】就假意推尊楚怀王为义帝，实际上不听从他的命令。\par
\noindent 【8.26.1】正月，项羽自立为西楚霸王，在梁、楚地区的九个郡称王，建都彭城。\par
\noindent 【8.26.2】背弃原来的约定，改立沛公为汉王，在巴、蜀、汉中称玉，建都南郑。\par
\noindent 【8.26.3】把关中瓜分为三，封立秦朝的三个将领：章邯为雍王，建都废丘：司马欣为塞王，建都栎阳；董翳为翟王，建都高奴。\par
\noindent 【8.26.4】封楚将瑕丘申阳为河南王，建都洛阳。\par
\noindent 【8.26.5】封赵将司马印为殷王，建都朝歌。\par
\noindent 【8.26.6】赵王歇迁徙代地称王。\par
  \item
\noindent 【8.20.13】回军攻打胡阳，遇到番君的别将梅鋗，与他一起，迫使析县、郦县投降。\par
\noindent 【8.20.14】派遣魏人宁昌出使秦关中，使者没有回来。\par
\noindent 【8.20.15】这时章邯已经带领全军在赵地投降项羽了。\par
\noindent 【8.21.1】起初，项羽和宋义北进援救赵，等到项羽杀死宋义，代替他为上将军，许多将领和黥布都从属项羽。\par
\noindent 【8.21.2】打垮了秦将王离的军队，使章邯投降，诸侯都归附了他。\par
\noindent 【8.21.3】等到赵高已经杀了秦二世，派人来见沛公，想要定约瓜分关中称王，沛公以为是诈骗，就采用张良的计策，派郦生、陆贾去游说秦军将领，用私利相诱，趁机袭击武关，攻破了关口。\par
\noindent 【8.21.4】又和秦军在蓝田南面交战，增设疑兵，多树旗帜，所经过的地方不许掳掠。\par
\noindent 【8.21.5】秦地的群众很高兴，秦军懈怠了，因此大破秦军。\par
\noindent 【8.21.6】又在蓝田北面接战，再次打败秦军。\par
\noindent 【8.21.7】乘胜追击，彻底打垮了秦军。\par
\noindent 【8.22.1】汉元年十月，沛公的军队先于各路诸侯到达霸上。\par
  \item
\noindent 【8.10.4】刘季用帛写了一封信，\par
\noindent 【8.10.5】射到城上，告诉沛县父老说：“天下苦于秦朝的暴政已经很久了。现在父老为沛令守城，但各国诸侯都已起事，（一旦城破，）就要屠戮沛县。如果沛县父老共同起来杀死沛令，选择子弟中可以立为首领的做领导，以响应诸侯军，那就能保全身家性命。不然的话，父子全遭杀害，死得毫无意义。”\par
\noindent 【8.10.6】父老们就率领子弟共同杀了沛令，打开城门，迎接刘季，想让他做柿县县令。\par
\noindent 【8.10.7】刘季说：“天下正在混乱当中，诸侯都已起事，如果推选的将领不胜任，就会一败涂地。我不是吝惜自己的生命，只怕才劣力薄，不能保全父兄子弟。这是件大事，希望另外共同推选一位能够胜任的人。”\par
\noindent 【8.10.8】萧何、曹参等都是文官，看重身家性命，怕事情不成，秦朝会诛灭他们的全族，所以都推让刘季。\par
\noindent 【8.10.9】父老们都说：“我们平时听到刘季许多奇异的事情，看来刘季是该显贵的。而且又经过占卜，没有比刘季更吉利的。”\par
\noindent 【8.10.10】这时刘季再三谦让，大家都不敢担任，最后还是立刘季为沛公。\par
\noindent 【8.10.11】在沛县衙门的庭院里祭祀黄帝和蚩尤，又用牲血衅鼓旗。\par
\noindent 【8.10.12】旗子一律红色，因为刘季所杀蛇是白帝的儿子，杀蛇的是赤帝的儿子，所以崇尚赤色。\par
\noindent 【8.10.13】于是少年子弟和有势的官吏，如萧何、曹参、樊哙等人，都为沛公征集兵员，集合了两三千人，攻打胡陵、方与，回军固守丰邑。\par
\noindent 【8.11.1】秦二世二年，陈胜将领周章的军队西至戏水而还。\par
  \item
\noindent 【8.21.7】乘胜追击，彻底打垮了秦军。\par
\noindent 【8.22.1】汉元年十月，沛公的军队先于各路诸侯到达霸上。\par
\noindent 【8.22.2】秦王于婴素车白马，用丝带系着脖子，封了皇帝的印玺和符节，在轵道旁投降。\par
\noindent 【8.22.3】将领们有的主张杀死秦王。\par
\noindent 【8.22.4】沛公说：“当初楚怀王派遣我，本来是因为我能宽大容人。况且人家已经降服，又杀死人家，不吉利。”\par
\noindent 【8.22.5】于是就把秦王交给了官吏，向西进入咸阳。\par
\noindent 【8.22.6】沛公想要留在宫殿中休息，樊哙、张良劝说后，才封闭了秦宫的贵重珍宝、财物和库房，回军霸上。\par
\noindent 【8.22.7】召集各县的父老、豪杰说：“父老们苦于秦朝的严刑峻法已经很久了，诽谤朝政的要灭族，相聚议论的要在街市上处斩。我和诸侯们约定，先入关的在关中称王，我应当称王关中。同父老们约定，法律只有三章：杀人的处死，伤人和抢劫的处以与所犯罪相当的刑罚。其余的秦朝法律全部废除。官吏和百姓都要安居如故。我所以到这里来，是为父老们除害，不会有欺凌暴虐的行为，不要害怕。我所以回军霸上，是等待诸侯们到来制定共同遵守的纪律。”\par
\noindent 【8.22.8】沛公派人与秦朝官吏巡行县城乡间，告谕百姓。\par
\noindent 【8.22.9】秦地的百姓大为高兴，争先恐后地拿出牛羊酒食款待士兵。\par
\noindent 【8.22.10】沛公又谦让不肯接受，说：“仓库的谷子很多，不缺乏，不愿破费百姓。”百姓更加高兴，唯恐沛公不做秦王。\par
\end{enumerate}
答案：\answerline{2.2cm}


\end{document}
